\scanpage{5}
équivaut \unsure{à|a} la condition
\[
j\mid G_{m\,\mathbb{R}}=i_{\mathbb{R}}\tag*{(a)}
\]

Pour la suite, $i$ est supposé donné, \deleted{\unsure{Considérons} \unsure{un} $j$ compatible avec $i$ au sens précédent.}\inserted{de $\overline{m}$ qu'un} objet inversible \inserted{$\mathbb{Q}(-1)$} de $\operatorname{Rep}(G)$ à isom près, i.e. un hom \unsure{$\varepsilon$}
\[
\varepsilon:G\longrightarrow G_m
\]
On a alors, pour un $j$ donné, que $\mathbb{Q}(1)$ est de bidegré $(1,1)$ si et seulement si on a les conditions
\[
\varepsilon j(\lambda)=\lambda^2\quad\text{si }\lambda\in\mathbb{R}^{*},\quad
\deleted{\unsure{j\mathbb{U}=1}}\quad
\varepsilon(j(u))=1\quad\text{si }u\in\mathbb{U}\tag*{(*)}
\]
i.e.
\[
\varepsilon j\deleted{\unsure{\mathbb{S}|j}}=\mathrm{Norme}:\unsure{\mathbb{S}|\mathrm{S}}\longrightarrow G_{m\,\mathbb{R}}\tag*{(b)}
\]
On suppose fixé $i$ et $\varepsilon$ satisfaisant \deleted{\unsure{$(a)$}}\inserted{$\varepsilon i(\lambda)=\lambda^2$}, et on s'intéresse aux $j$ satisfaisant a) et b).

\inserted{Soit $j$ donné.} Cherchons, pour $V\in\operatorname{Ob}\operatorname{Rep}(G)$ donné\inserted{ de degré $p$}, \unsure{d'où|D'où} structure de Hodge sur $V_{\mathbb{C}}$, (ou sur $V_{\mathbb{R}}$) quelles sont les \deleted{polarisations de cette structure de Hodge, provenant}\inserted{formes de $R(G)$}
\[
V\times V\xrightarrow{\varphi}\mathbb{Q}(-p)
\]
qui définissent une \uline{polarisation} des structures de Hodge correspondantes. Ce sont celles qui satisfont \unsure{sur|une} la condition que
\[
\varphi_{\mathbb{R}}(x,Cy)\quad\text{forme \inserted{\uline{symétrique}} \uline{définie positive sur} } V_{\mathbb{R}}\times V_{\mathbb{R}}\tag*{(**)}
\]
% archive sheet number: 4

